\documentclass{article}
\usepackage{amsmath}
\usepackage[a4paper,margin=1in]{geometry}

\begin{document}

\title{Data Science Interview Puzzles -- Questions and Answers}
\date{}
\maketitle

\section*{Puzzle 1: Piano Tuners in Hyderabad}

\textbf{Question:} Estimate the number of piano tuners in Hyderabad.

\textbf{Answer:}
This is a Fermi estimation problem. A reasonable approach is:
\begin{enumerate}
    \item Estimate the number of pianos in Hyderabad.
    \item Estimate the average number of tunings per piano per year.
    \item Estimate how many tunings a single tuner can perform annually.
    \item Compute
    \[
    \text{Tuners} = \frac{\text{Total Tunings Per Year}}{\text{Tunings Per Tuner Per Year}}.
    \]
\end{enumerate}
The exact number is less important than a logical decomposition and reasonable assumptions.

\section*{Puzzle 2: Model Has 99.9\% Accuracy}

\textbf{Question:} A model has 99.9\% accuracy on 1 million rows. Is it good?

\textbf{Answer:}
The first question should be:

\textit{``What is the class distribution?''}

If the dataset is highly imbalanced, accuracy can be misleading.
For example, if 99\% of observations belong to one class, a model predicting only the majority class achieves 99\% accuracy while being useless.

Additional metrics to inspect:
\begin{itemize}
    \item Precision
    \item Recall
    \item F1-score
    \item ROC-AUC or PR-AUC
\end{itemize}

\section*{Puzzle 3: Recommendation Algorithms}

\textbf{Question:}
Algorithm A has CTR = 5.1\%. Algorithm B has CTR = 5.4\%.
Can we conclude B is better?

\textbf{Answer:}
No.

Things to check:
\begin{itemize}
    \item Sample size
    \item Randomized A/B test design
    \item Statistical significance
    \item Confidence interval of the lift
    \item Business significance
\end{itemize}

A small increase may simply be random variation.

\section*{Puzzle 4: Expected Tosses Until First Head}

\textbf{Question:}
A fair coin is tossed until the first Head appears.
What is the expected number of tosses?

\textbf{Answer:}
Let the expectation be $E$.

\[
E=\frac12(1)+\frac12(1+E)
\]

Solving:

\[
E=2.
\]

\[
\boxed{E=2}
\]

This is an example of \textbf{first-step analysis}.

\section*{Puzzle 5: Feature X Increased Sales}

\textbf{Question:}
A manager says:

\textit{``Sales increased after Feature X launched, therefore Feature X caused the increase.''}

\textbf{Answer:}
Correlation is not causation.

Potential confounders:
\begin{itemize}
    \item Seasonality
    \item Marketing campaigns
    \item Existing business trend
    \item Competitor actions
    \item Holidays or special events
\end{itemize}

A better approach is an A/B test or a causal inference framework.

\section*{Puzzle 6: Traffic Lights in Hyderabad}

\textbf{Question:}
Estimate the number of traffic lights in Hyderabad.

\textbf{Answer:}
One possible decomposition:
\begin{enumerate}
    \item Estimate the number of major intersections.
    \item Estimate the proportion that use traffic lights.
    \item Estimate the average number of signal units per intersection.
\end{enumerate}

The focus is on assumptions and reasoning rather than the final numerical answer.

\section*{Puzzle 7: Cancer Detection Metrics}

\textbf{Question:}
For a cancer detection model, which is more important: Accuracy, Precision, or Recall?

\textbf{Answer:}
Usually Recall is most important because false negatives are very costly.

\[
\text{Recall}=\frac{TP}{TP+FN}
\]

Missing a cancer patient is often worse than generating an extra false alarm.

However, the correct metric depends on the business and medical costs associated with false positives and false negatives.

\section*{Puzzle 8: Heavy Ball Puzzle}

\textbf{Question:}
One of eight balls is heavier than the others.
Using a balance scale, what is the minimum number of weighings required to identify it?

\textbf{Answer:}
Two weighings.

Reason:

Each weighing has three outcomes:
\begin{itemize}
    \item Left heavier
    \item Right heavier
    \item Equal
\end{itemize}

Two weighings provide
\[
3^2=9
\]
possible outcome patterns, which is enough to distinguish among eight candidate balls.

\section*{Puzzle 9: First Look at a CSV Dataset}

\textbf{Question:}
You receive a CSV containing 100 columns.
What is the first thing you do?

\textbf{Answer:}
First understand the schema:
\begin{itemize}
    \item What is the target variable?
    \item Classification or regression?
    \item Data types?
    \item Missing values?
    \item Duplicates?
\end{itemize}

Then examine the target distribution and begin exploratory data analysis.

\end{document}
